
\subsubsection{5.1 Sample Efficiency (H-E1)}

At N=1,000 training models:

\begin{table}[htbp]
\centering
\resizebox{\columnwidth}{!}{%
\begin{tabular}{ll}
\toprule
Model & $R^2$ \\
\midrule
NFN & 0.9524 \\
MLP-Matched & 0.3529 \\
**Difference** & **0.5995** \\
\bottomrule
\end{tabular}
}
\end{table}

The $R^2$ difference of 0.60 exceeds the 0.05 threshold by a factor of 12. NFN shows minimal overfitting (train loss 0.0024, validation loss 0.0029), while MLP shows severe overfitting (train loss ~0, validation loss 0.045).

\subsubsection{5.2 NFN Equivariance Verification (H-M1, H-M2)}

Applying 10 random neuron permutations to test MLP weights and measuring NFN output:

\begin{table}[htbp]
\centering
\resizebox{\columnwidth}{!}{%
\begin{tabular}{ll}
\toprule
Metric & Value \\
\midrule
Mean prediction & 0.8760 \\
Standard deviation & 8.$22\times 10^{-8}$ \\
Max deviation & 1.$19\times 10^{-7}$ \\
Invariance correlation & 0.99999986 \\
\bottomrule
\end{tabular}
}
\end{table}

All 11 predictions (original plus 10 permuted) differ only at floating-point precision. This confirms NFN's mathematical equivariance property.

\subsubsection{5.3 Untrained MLP Variance (H-M3)}

An untrained (randomly initialized) MLP produces highly variable outputs under the same permutations:

\begin{table}[htbp]
\centering
\resizebox{\columnwidth}{!}{%
\begin{tabular}{ll}
\toprule
Metric & Value \\
\midrule
Mean prediction & -0.0607 \\
Standard deviation & 0.0118 \\
Coefficient of variation & 0.194 \\
Max deviation & 0.0229 \\
\bottomrule
\end{tabular}
}
\end{table}

The coefficient of variation (0.194) exceeds the 0.1 threshold, confirming that MLPs have no built-in permutation invariance.

\subsubsection{5.4 MLP at Large Scale (H-M5)}

At N=40,547 training models (3 seeds):

\begin{table}[htbp]
\centering
\resizebox{\columnwidth}{!}{%
\begin{tabular}{llll}
\toprule
Metric & Mean & Std & 95\% CI \\
\midrule
Test $R^2$ & 0.9860 & 0.0003 & [0.985, 0.987] \\
Invariance & 0.6265 & 0.0103 & [0.595, 0.658] \\
\bottomrule
\end{tabular}
}
\end{table}

MLP achieves near-perfect prediction accuracy but fails the invariance criterion (0.63 < 0.80). This is the central negative result: data diversity alone does not teach permutation invariance.

\subsubsection{5.5 Comparison Across Scales}

\begin{table}[htbp]
\centering
\resizebox{\columnwidth}{!}{%
\begin{tabular}{lll}
\toprule
Condition & $R^2$ & Invariance \\
\midrule
MLP N=1K & 0.004 & 0.64 \\
MLP N=40K & 0.986 & 0.63 \\
NFN (any N) & ~0.95 & 1.0 \\
\bottomrule
\end{tabular}
}
\end{table}

MLP invariance does not increase meaningfully from N=1K to N=40K despite a dramatic increase in predictive performance.

\subsubsection{5.6 Summary of Hypothesis Outcomes}

\begin{table}[htbp]
\centering
\resizebox{\columnwidth}{!}{%
\begin{tabular}{lll}
\toprule
Hypothesis & Result & Status \\
\midrule
H-E1 & $\Delta =0.60$ > 0.05 & **PASSED** \\
H-M1 & 1.$19\times 10^{-7}$ < $10^{-5}$ & **PASSED** \\
H-M2 & 0.9999999 > 0.99 & **PASSED** \\
H-M3 & CV=0.194 > 0.1 & **PASSED** \\
H-M4 & 0.64 > 0.5 & FAILED \\
H-M5 & 0.63 < 0.8 & **FAILED** \\
\bottomrule
\end{tabular}
}
\end{table}

The H-M5 failure is the most informative result: it falsifies the hypothesis that MLPs learn invariance from data diversity at scale.

